\section{Continuous envelopes for the finite certificates}\label{ce:section}

The finite calculations use two envelopes. The first freezes a
moment-comparison law on an entire parameter box. The second keeps the
same noise in the comparison variance and moment cost until a derivative
test licenses its elimination. The tests below are obligations of each
certificate, rather than consequences asserted for every parameter choice.

\begin{proposition}[A fixed-reference box]\label{ce:fixed-box}
Let $p\in[p_l,p_u]\subset(0,1/2)$ and $\ell\le\lambda\le u$, with
$\nu=np_l(1-p_l)+\ell>0$. Suppose a reference expression $A(p,t)$,
independent of the actual noise, and a differentiable moment floor
$g(p,\lambda)$ give
\[
 R(t)\le A(p,t)-\Phi\left(\frac{k+t-np}{\sqrt{npq+\lambda}}\right)
                      -c_ng(p,\lambda).
\]
If $g_\lambda\ge b$ throughout the box, then
\begin{align}\label{ce:fixed-box-bound}
 R(t)\le{}&A(p,t)-\Phi\left(\frac{k+t-np}{\sqrt{npq+\ell}}\right)
                      -c_ng(p,\ell)+J(t),\\
 J(t)={}&(u-\ell)
 \left\{\frac{(k+t-np_l)_+\phi(0)}{2\nu^{3/2}}-c_nb\right\}_+.
 \notag
\end{align}
If $A(p,\cdot)$ is convex, the upper function has the phase
semiconvexity constant in \eqref{pp:gaussian-curvature} with
$\lambda_l=\ell$.
\end{proposition}
\begin{proof}
For $\mu=k+t-np$, differentiation of the explicit surrogate gives
\[
 \frac\partial{\partial\lambda}
 \left\{-\Phi\left(\frac\mu{\sqrt{npq+\lambda}}\right)-c_ng(p,\lambda)\right\}
 =\frac{\mu\phi(\mu/\sqrt{npq+\lambda})}{2(npq+\lambda)^{3/2}}
                           -c_ng_\lambda.
\]
Its upper bound is the expression in braces defining $J$.
Integrate from $\ell$ to the actual $\lambda$. The nonnegative
correction $J$ is convex in $t$, so only the Gaussian term contributes
to the lower phase-curvature bound.
\end{proof}

Here $A$ may be the current-count price
\[
 A(p,t)=F_k+\eta(p)\{t+Q_Z(t+\pi_k/\eta(p))\},
\]
or a fixed reflected-square reference price. For example, upper bounds
$a,z_0,z_1$ on the first group's inward support and the two conditional
variances give
\[
 Z=\sum_{i=1}^{n-k}T_{a,z_0,i}+\sum_{i=1}^{k}T_{O_1,z_1,i}.
\]
Lemma~\ref{mc:two-moment} makes this one law valid throughout the box.
The necessary resource plane \eqref{mc:signed-plane} is retained when
bounding $\lambda(1-\alpha)$ and $\lambda\alpha$ for $z_0,z_1$.
An inward cap and an allocation lower bound similarly give
\[
 a_0\le\delta-\frac{\alpha\lambda}{npO_1}.
\]
These relations keep the variance and support corners within a single
proved resource budget.

The probability direction can also be bounded without a probability
mesh. At a fixed phase, put $a(p)=t+\pi_k(p)/\eta(p)$.
For the term $\eta Q_Z(a)$, a lower second-derivative bound is supplied by
\begin{equation}\label{ce:probability-curvature}
 \eta''Q_Z(a)+2\eta'Q_Z'(a)a'+\eta Q_Z'(a)a''.
\end{equation}
The omitted term is $\eta Q_Z''(a)(a')^2\ge0$, interpreted weakly
at an activation point. The remaining binomial, Gaussian and moment-floor
derivatives are explicit. Bounds on these expressions and on
$-1\le Q_Z'\le0$ give a uniform probability chord by
Lemma~\ref{pp:chord}. Face bounds obtained from different moment
estimates are combined through the same fixed-reference envelope with
the actual law's original moment and variance.

\begin{proposition}[A coherent reflected-square envelope]\label{ce:coherent}
Let $j=k+1<n$ and fix a legal square width $h$ for reference count $j$.
Suppose the centered reflected residual summands have one common upper
bound $b>0$. Set
\[
 V_j=\frac\lambda n\left\{\frac{(n-j)(1-\alpha)}q+\frac{j\alpha}p\right\},
 \qquad K(V,y)=\E(Z_b(V)-y)_+^2.
\]
Let $G(p,\lambda,\alpha)$ be a valid actual-moment floor, convex in
$\alpha\in[0,1]$ and differentiable there. A positive affine inward
cap in \eqref{mc:allocation-floor} supplies such a floor.
Define
\begin{align}\label{ce:coherent-price}
 E(p,\lambda,\alpha,t)={}&F_k+
     \frac{\pi_j}{h^2}K(V_j,1-t-h)
 -\Phi\left(\frac{k+t-np}{\sqrt{npq+\lambda}}\right)
 -c_nG(p,\lambda,\alpha).
\end{align}
Suppose, on the whole stated probability, noise and phase domain,
\begin{equation}\label{ce:alpha-gate}
 c_nG_\alpha(p,\lambda,0)
 \ge\frac{\pi_j\lambda(j-np)_+}{npq\,h^2},
\end{equation}
and
\begin{align}\label{ce:noise-gate}
 \frac{\pi_j(n-j)}{nq h^2}K_V(V_j(p,\lambda,0),1-t-h)
 +\frac{(k+t-np)\phi((k+t-np)/\sqrt{npq+\lambda})}
                  {2(npq+\lambda)^{3/2}}\notag\\
  {}-c_nG_\lambda(p,\lambda,0)\ge0.
\end{align}
Then for every actual $\lambda\in[\ell,u]$ and allocation,
\begin{equation}\label{ce:coherent-bound}
 R(t)\le E(p,\lambda,\alpha,t)\le E(p,u,0,t).
\end{equation}
The zero-noise boundary is included by continuity.
\end{proposition}
\begin{proof}
Apply \eqref{mc:clipped-square} at count $j$ and phase $t-1$,
then apply \eqref{mc:bentkus} to the reflected residual sum.
This proves the first upper bound with the same actual third moment,
replaced only by its valid lower floor $G$.

By \eqref{mc:variance-derivative}, $0\le K_V\le1$, and
$\partial_\alpha V_j=\lambda(j-np)/(npq)$.
Convexity gives $G_\alpha(p,\lambda,\alpha)\ge G_\alpha(p,\lambda,0)$.
Thus \eqref{ce:alpha-gate} implies $E_\alpha\le0$.
On the formal zero-allocation face, \eqref{ce:noise-gate} is exactly
$\partial_\lambda E(p,\lambda,0,t)\ge0$. Integrating the two
monotonicities proves \eqref{ce:coherent-bound}.
\end{proof}

Zero allocation here is an endpoint of the majorizing function; it need
not describe an admissible original law. The first elimination uses the
universal bound $0\le K_V\le1$. Only afterwards is a one-clip formula
needed to evaluate the formal face or obtain the sharper lower derivative
in \eqref{ce:noise-gate}. Its next-atom condition is checked on that face.
The resulting explicit function $E(p,u,0,t)$ is bounded in probability
and phase by the weak-curvature rules of Section~\ref{pp:section}.

\begin{corollary}[An allocation interval]\label{ce:allocation-interval}
In the same reflected-square setting, retain one allocation interval
$[\alpha_-,\alpha_+]$. Suppose the moment floor $G$ is nondecreasing
in allocation throughout the domain. For either endpoint
$\xi\in\{\alpha_-,\alpha_+\}$ define
\begin{align}\label{ce:two-corners}
 E_{\xi}(p,\lambda,t):={}&F_k+
 \frac{\pi_j}{h^2}K(V_j(p,\lambda,\xi),1-t-h)\notag\\
 &-\Phi\left(\frac{k+t-np}{\sqrt{npq+\lambda}}\right)
 -c_nG(p,\lambda,\alpha_-).
\end{align}
Then, throughout the original probability interval, including a crossing
of $j=np$,
\begin{equation}\label{ce:variance-endpoint-max}
 R(t)\le\max\{E_{\alpha_-}(p,\lambda,t),
              E_{\alpha_+}(p,\lambda,t)\}.
\end{equation}
If $j\ge np$ everywhere, only $\xi=\alpha_+$ is needed;
if $j\le np$ everywhere, only $\xi=\alpha_-$ is needed.
If $\partial_\lambda E_\xi\ge0$ on $[\ell,u]$ for every retained
branch, replace $\lambda$ by $u$ in the resulting maximum.
\end{corollary}
\begin{proof}
The variance $V_j$ is affine in allocation, so its maximum on an
interval occurs at an endpoint. Since $K$ increases with variance
and the moment cost is at least $G(p,\lambda,\alpha_-)$,
\eqref{ce:variance-endpoint-max} follows. The derivative
$\partial_\alpha V_j=\lambda(j-np)/(npq)$ selects a single endpoint
when its sign is fixed. Each retained explicit branch keeps the
same actual $\lambda$ and can be increased to $u$ under its own
noise derivative condition. Zero noise is included by continuity.
\end{proof}

The noise gate in this corollary is explicit. With
$\nu_\xi=\{(n-j)(1-\xi)/q+j\xi/p\}/n$, it is
\[
 \frac{\pi_j\nu_\xi}{h^2}K_V(\nu_\xi\lambda,1-t-h)
 +\frac{(k+t-np)\phi((k+t-np)/\sqrt{npq+\lambda})}
               {2(npq+\lambda)^{3/2}}
 -c_nG_\lambda(p,\lambda,\alpha_-)\ge0.
\]
Probability and phase bounds are also applied to each explicit branch
before taking their maximum. A single branch across a crossing is
only an auxiliary expression; the maximum in
\eqref{ce:variance-endpoint-max} gives the actual residual bound.
Both allocation rules use one endpoint test. The affine-cap moment
floor is convex by Section~\ref{mc:section}, so
$G_\alpha(p,\lambda,\alpha_-)\ge0$ implies $G(\alpha)\ge G(\alpha_-)$.
The stronger endpoint payment in \eqref{ce:alpha-gate}, with $0$
replaced by $\alpha_-$, instead makes the full price nonincreasing.
To evaluate that endpoint, write $\delta=A-\lambda(H+r\alpha)>0$,
$B=(1-\alpha)/\sqrt q+\alpha/\sqrt p$ and
$B_\alpha=1/\sqrt p-1/\sqrt q$. Direct differentiation gives
\begin{equation}\label{ce:signed-allocation-derivative}
 G_\alpha=\frac{s}{(v+s)^{3/2}}
 \left\{3(q-p)+\lambda r+s\left(
 \frac{2BB_\alpha}{\delta}+\frac{\lambda rB^2}{\delta^2}\right)\right\}.
\end{equation}
For $p\in[p_l,p_u]\subset(0,1/2)$, $0\le\lambda\le L$,
$\delta\ge\delta_*>0$ and $B\le B_*$, a lower bound on the
braced quantity is
\begin{equation}\label{ce:signed-allocation-gate}
 3(1-2p_u)+L\min(0,r)
       +\frac{L^2\min(0,r)B_*^2}{n\delta_*^2}.
\end{equation}
This is an optional sufficient bound for the endpoint test.
The omitted $2BB_\alpha$ term is nonnegative, so a positive bound
licenses the choice for either sign of $r$.
Zero noise again follows by continuity.
